\documentclass[10pt,a4paper]{article}
\usepackage[margin=2cm]{geometry}
\usepackage{booktabs,float,amsmath,array}
\usepackage[hidelinks]{hyperref}
\hypersetup{pdftitle={CS433 Assignment 3 Report, Group 1},pdfauthor={Ashik Stenny, Spandan Pati, Vedhanth Balasubramanian}}
\setlength{\parskip}{4pt}\setlength{\parindent}{0pt}
\IfFileExists{gen/macros.tex}{\newcommand{\devname}{NVIDIA GeForce RTX 3050 Laptop GPU}
\newcommand{\devcomputecapability}{8.6}
\newcommand{\devsms}{16}
\newcommand{\devmaxthreadsperblock}{1024}
\newcommand{\devmaxthreadspersm}{1536}
\newcommand{\devmaxblockspersm}{16}
\newcommand{\devmaxresidentthreads}{24576}
\newcommand{\devwarpsize}{32}
\newcommand{\devglobalmemmb}{3770}
\newcommand{\devsharedmemperblockkb}{48}
\newcommand{\devsharedmempersmkb}{100}
\newcommand{\devlcachekb}{1536}
\newcommand{\devmembuswidthbits}{128}
\newcommand{\devmemclockmhz}{5501}
\newcommand{\devpeakbandwidthgbs}{176.0}
\newcommand{\ponestaticms}{1464}
\newcommand{\ponecyclicms}{1461}
\newcommand{\ponedynamicms}{1370}
\newcommand{\ponechunk}{1}
\newcommand{\ponechunkthreads}{$2^{20}$}
\newcommand{\ponebs}{1024}
\newcommand{\ponebestthreads}{$2^{22}$}
\newcommand{\ponebestms}{1370}
\newcommand{\ptwogroup}{32}
\newcommand{\ptwobs}{128}
\newcommand{\ptwobestthreads}{$2^{26}$}
\newcommand{\ptwobestms}{351.0}
\newcommand{\qonedynspeedupmin}{5.31}
\newcommand{\qonedynspeedupmax}{8.63}
\newcommand{\qonedynvsstaticmin}{1.09}
\newcommand{\qonedynvsstaticmax}{1.75}
\newcommand{\qonespeedupmin}{3.93}
\newcommand{\qonespeedupmax}{5.77}
\newcommand{\qoneallverified}{All}
\newcommand{\ptwodynchunk}{1}
\newcommand{\ptwodynbestms}{259.5}
\newcommand{\ptwodynbestthreads}{$2^{18}$}
\newcommand{\ptwodyngroup}{32}
\newcommand{\qoneseedratioonemin}{3.3}
\newcommand{\qoneseedratioonemax}{5.7}
\newcommand{\qoneseedratiotwomin}{1.27}
\newcommand{\qoneseedratiotwomax}{1.70}
\newcommand{\qoneseedsameorder}{15}
\newcommand{\qtwobasicbestt}{$2^{14}$}
\newcommand{\qtwobasicbestbs}{32}
\newcommand{\qtwobasicbestms}{16.60}
\newcommand{\qtwovonems}{16.62}
\newcommand{\qtwovonegbps}{64.61}
\newcommand{\qtwovtwoms}{6.630}
\newcommand{\qtwovtwogbps}{161.96}
\newcommand{\qtwovthreems}{6.449}
\newcommand{\qtwovthreegbps}{166.53}
\newcommand{\qtwovfourms}{6.349}
\newcommand{\qtwovfourgbps}{169.13}
\newcommand{\qtwovfivems}{6.269}
\newcommand{\qtwovfivegbps}{171.30}
\newcommand{\qtwospeedupmin}{1.0}
\newcommand{\qtwospeedupmax}{4.7}
}{}
\newcommand{\inputif}[1]{\IfFileExists{gen/#1}{\input{gen/#1}}{\textit{[table #1 not generated yet]}}}

\title{CS433 Assignment 3: Report\\\large Group \#1}
\author{Ashik Stenny \and Spandan Pati \and Vedhanth Balasubramanian}
\date{November 2026}
\begin{document}
\maketitle

%=====================================================================
\section{Experimental setup}
We ran all experiments on one of our laptops with the charger connected. We
noticed that on battery the GPU runs at roughly half its clock, so we did not use
any runs made on battery.

\begin{table}[H]\centering\small
\caption{Hardware and software.}\label{tab:setup}
\begin{tabular}{ll}\toprule
GPU & \devname{} (compute capability \devcomputecapability) \\
SMs & \devsms{}, max \devmaxthreadspersm{} threads and \devmaxblockspersm{} blocks per SM \\
Max resident threads & \devmaxresidentthreads{} (= \devsms{} $\times$ \devmaxthreadspersm) \\
Memory & \devglobalmemmb{} MB GDDR6, \devmembuswidthbits-bit bus, peak \devpeakbandwidthgbs{} GB/s \\
L2 cache / shared memory & \devlcachekb{} KB / \devsharedmemperblockkb{} KB per block \\
CPU & AMD Ryzen 7 5800HS, 8 cores / 16 threads \\
Software & CUDA 12.4 (\texttt{nvcc -O3 -arch=sm\_86}), gcc \texttt{-O3 -march=native -fopenmp} \\
\bottomrule\end{tabular}\end{table}

GPU kernels in Problem~1 are timed with CUDA events; Problem~2 uses
\texttt{gettimeofday} around launch + \texttt{cudaDeviceSynchronize}/copy of the
result, so the time includes everything needed to get the final scalar on the
host. Every time in the final comparison tables and the best-configuration tables
is the average of 3--5 repetitions; only the two large exploratory sweeps (P1
chunk size and P2 group size versus thread count) use one repetition per point.
Repeated runs of the same configuration are not identical. In the final comparison
table the minimum and the average of the repetitions differ by about 2\% on average
and by at most about 10\%. The spread is largest (up to about 13\%) for P1 with static
assignment and few threads, where the run time depends on which threads happen to get
the expensive late inserts, and that changes from run to run. These differences are
much smaller than the differences we compare (mostly factors of 1.3 to 15), so we
think they do not change our conclusions. Before and after every series of runs we also ran a short bandwidth
test to check that the GPU was at full clock (about 170~GB/s). We added this check
after we found that some runs made after a suspend/resume of the laptop were
throttled; those runs were discarded and repeated.

%=====================================================================
\section{How to compile and run}
We submit three source files for Problem~1 (\texttt{p1.cu}, \texttt{p2.cu},
\texttt{gen\_ops.c}) and six for Problem~2 (\texttt{q2\_v1\_basic.cu} to
\texttt{q2\_v5\_opt.cu} and \texttt{q2\_omp.c}). Each one is self-contained (the helper
code shared by the programs of a problem is pasted into each file), so it compiles on
its own and no Makefile is needed. The scripts we used to run the experiments and to
make the tables are not part of the submission.
\texttt{-arch=native} (CUDA $\ge$ 11.6) compiles for the GPU of the machine; with
an older \texttt{nvcc} use the GPU's architecture instead, e.g.\ \texttt{-arch=sm\_61}
(GTX 1080) or \texttt{-arch=sm\_86} (A40, RTX 30xx).

\textbf{Problem 1} (directory \texttt{q1/}):
{\footnotesize\begin{verbatim}
gcc  -O2 -o gen_ops gen_ops.c
nvcc -O3 -arch=native -std=c++14 -o p1 p1.cu
nvcc -O3 -arch=native -std=c++14 -o p2 p2.cu

./gen_ops <opfile> <nops> <ins%> <del%> <srch%> [N=2^30] [seed=12345]
./p1 <opfile> <nops> <nthreads> [blocksize=1024] [static|cyclic|dynamic] [chunk=1]
./p2 <opfile> <nops> <nthreads> [blocksize=128] [groupsize=32] [static|dynamic] [chunk=1]

# example: 2^22 operations, 80% insert / 20% search, both programs at their best settings
./gen_ops ops.bin 4194304 80 0 20
./p1 ops.bin 4194304 4194304 1024 dynamic 1
./p2 ops.bin 4194304 67108864 128 32 dynamic 1
\end{verbatim}}
The first three arguments of \texttt{p1}/\texttt{p2} are the ones the assignment
prescribes (operation file, number of operations, thread count); the optional ones
default to the best values found in Section~\ref{sec:q1res}, except that \texttt{p2}
uses the static (cyclic) group schedule unless \texttt{dynamic} is given; the best
P2 configuration is \texttt{128 32 dynamic 1}, as in the example. Each run prints the
time (average and minimum over the repetitions) and \texttt{verify=OK} once every
return value and the final table have been checked (paragraph \emph{Verification} in Section~\ref{sec:q1alg}).

\textbf{Problem 2} (directory \texttt{q2/}; $m$, $n$, $t$ powers of two):
{\footnotesize\begin{verbatim}
nvcc -O3 -arch=native -Xcompiler -fopenmp -lgomp -o q2_v1_basic q2_v1_basic.cu
#   ... the same for q2_v2_coalesced, q2_v3_tree, q2_v4_shmem, q2_v5_opt
gcc  -O3 -march=native -fopenmp -o q2_omp q2_omp.c -lm

./q2_v1_basic <m> <n> <t> [blocksize]     # contiguous chunk per thread
./q2_v2_coalesced ...                     # + memory access coalescing
./q2_v3_tree ...                          # + tree reduction
./q2_v4_shmem ...                         # + shared memory
./q2_v5_opt ...                           # + float4 loads, ILP, warp shuffle (best)
./q2_omp <m> <n> <t>                      # OpenMP version used for the comparison
./q2_v5_opt 16384 16384 65536             # example
\end{verbatim}}
The environment variable \texttt{REPS} sets the number of timed repetitions
(default 5 for Problem~1, 3 for Problem~2).

%=====================================================================
\section{Problem 1: concurrent open-addressed hash table}

\subsection{Operation lists}
\texttt{gen\_ops} writes a binary file of \texttt{(op, key)} pairs of 32-bit
integers. Insert keys are \texttt{random() \% N} with $N=2^{30}$, which is
$64\times$--$1024\times$ larger than the table sizes used ($2^{20}$--$2^{24}$).
The list is segregated: all inserts, then all deletes, then all searches.
Delete keys are \emph{distinct} keys that were inserted (a second delete of
the same key would return 0); search keys are drawn from inserted keys that
are not deleted. Thus every delete and search returns 1. For the
50\%/50\% insert/delete mix every inserted key must be distinct, so the
generator redraws the (few) duplicate insert keys until there are enough
distinct keys. A redrawn key can itself collide with an existing key (probability
about $n/N$, i.e.\ 0.1--1.6\% here), so each round of redraws is followed by an
exact recount of the distinct keys, and further rounds run until the count is
sufficient (2--3 rounds in practice). This is also why the generator rejects mixes
with more deletes than possible keys ($n_d > N$), for which it could never finish.
The table size $n$ is computed by the programs as the number
of distinct insert keys (sorting on the host, not timed), and $h(k)=k \bmod n$.

\subsection{Hash table algorithm and correctness}\label{sec:q1alg}
The table is an \texttt{int} array with two reserved values,
\texttt{EMPTY}$=-1$ (set by \texttt{cudaMemset(0xFF)}) and
\texttt{TOMB}$=-2$ (a deleted key). All operations are lock-free:
\begin{itemize}
\item \textbf{insert($k$)}: probe from $h(k)$; if a slot holds $k$, return 0.
  On an \texttt{EMPTY} slot do \texttt{atomicCAS(EMPTY}$\to k$): if it
  succeeds, return 1; if it returns $k$, another thread has just inserted $k$
  (return 0); otherwise another key won the slot and probing continues.
  \textbf{Tombstones are reused:} the insert remembers the \emph{first} tombstone
  on its path but keeps scanning until it has proven that $k$ is absent (it reaches
  an \texttt{EMPTY} slot or has visited every slot), and only then claims that
  tombstone with \texttt{atomicCAS(TOMB}$\to k$); if another key took it first,
  the scan restarts.
\item \textbf{delete($k$)}: probe until $k$ (then \texttt{atomicCAS}($k\to$\texttt{TOMB}),
  return whether it succeeded) or \texttt{EMPTY} (return 0).
\item \textbf{search($k$)}: probe until $k$ (1) or \texttt{EMPTY} (0).
\end{itemize}
A deleted slot becomes a tombstone rather than \texttt{EMPTY} so that the probe
chains of other keys are not broken (a search for a key further along the chain
must not stop at the deleted slot). We argue that a key can never end up in two slots: all
inserters of $k$ walk the same probe sequence, the key slots they pass cannot
change during an insert kernel, and they therefore all compete with a CAS for
the same first free slot (the first tombstone, or the first \texttt{EMPTY} slot).
The scan before reusing a tombstone is needed because $k$ may already be
stored \emph{after} the tombstone. Within one kernel a slot only moves forward
(\texttt{EMPTY}/\texttt{TOMB}$\to k$ in an insert kernel, $k\to$\texttt{TOMB} in
a delete kernel), so as far as we can see a stale (e.g.\ L1-cached) read is
harmless: the CAS that follows always sees the true value. Because $n$ equals the number of distinct
keys, an insert of an absent key always finds a free slot: at most $n-1$ slots
hold other keys. We use a CAS on the slot itself instead of a lock
(\texttt{while(atomicCAS(\&lock,0,1));}), which would serialise the operations
and, as discussed in lecture~6, can deadlock inside a warp on pre-Volta GPUs. We did
not use \texttt{\_\_threadfence}: the only shared writes are single atomic CAS
operations, and no thread depends on the ordering of two different writes, so we do
not think a fence is needed.

\textbf{Design decision: segment-wise execution.} The host splits the
operation list into runs of identical operations and launches one kernel per
run. A kernel boundary is a global barrier, so this gives \emph{sequential
semantics} between operation types: every operation sees the effects of all
earlier operations of a different type (e.g.\ a search always sees every insert
before it). The expected result of every operation is then well defined and
can be verified exactly. A segregated list (as required) has at most three runs,
i.e.\ three kernel launches ($\approx$20~$\mu$s in total). A list that is not
segregated is also processed correctly, one kernel per run (the program prints a
warning), but every change of operation type costs a kernel launch of
5--10~$\mu$s, so a fully mixed list would be dominated by launch overhead.
Executing a mixed list in a single kernel would let an insert and a delete or
search of the same key run concurrently. The outcome would then depend on timing
(a search may or may not see a concurrent insert), the ``always returns 1''
property of positive queries would no longer hold, and concurrent inserts and
deletes of the same key would need a separate correctness argument. Since the
assignment requires segregated lists, we did not implement that mode.

\textbf{Verification} (after \emph{every} timed repetition, not timed). Before
the runs, the host replays the operation list segment by segment with a sorted
set of keys and derives the exact expected outcome for every key of every
segment. Inside one segment all operations run concurrently, so the rule is
per key: an insert segment must return 1 for exactly one occurrence of each
key that was absent (and 0 for all others), a delete segment must return 1 for
exactly one occurrence of each key that was present, and every search must
return whether its key was present. Only \emph{which} of several duplicate
operations wins is left open, since that depends on timing. After every run we
check
(i) every key group's return values against this plan,
(ii) that the table contains \emph{exactly} the expected set of keys (none
missing, none extra, none twice), and
(iii) that every key is reachable from its home slot (no \texttt{EMPTY} slot in
between, checked in one linear pass).
This check is stricter than comparing totals (two wrong results of different keys
cannot cancel out), and it also works for non-segregated lists. When we first
introduced it, it found a bug in our earlier version: tombstones were never
reused, so in a list that deletes and then re-inserts a key into a full table, the
re-insert found no free slot. This cannot happen with segregated lists (all
inserts precede all deletes), so no measured result was affected, but it led to
the tombstone reuse described above. All runs reported below passed
(column ``verified'' in Table~\ref{tab:q1final}). The programs were also checked
on hand-made mixed lists (fill the table, delete half the keys, re-insert them,
further deletes, inserts and searches) with every group size.

\textbf{Load factor.} The table is sized to the number of distinct keys, so
after the inserts it is 100\% full. With linear probing the expected cost of
filling a table completely grows like $n^{3/2}$ (the last insertions must
walk huge clusters), so the run time grows much faster than the number of
operations, and it is dominated by the last few percent of inserts. This
alone predicts $4^{3/2}=8\times$ the time for $4\times$ the operations, which fits
the step from $2^{20}$ to $2^{22}$ (9.4$\times$ for P1), but not the step from
$2^{22}$ to $2^{24}$ (48$\times$). We think at least part of the rest comes from the
modulo bias of $h(k)$ discussed in Section~\ref{sec:q1res}.

\subsection{P1: one thread per operation}
Thread count $T$ = blocks $\times$ block size (command-line arguments). Three
task assignments are implemented:
\textbf{static} (block): thread $i$ gets the contiguous ops
$[i\lceil c/T\rceil, (i+1)\lceil c/T\rceil)$;
\textbf{cyclic} (static round-robin): thread $i$ gets ops $i, i+T, \dots$, so a
warp reads 32 consecutive operation records (coalesced);
\textbf{dynamic}: a global counter per kernel; each thread grabs a chunk of
$C$ operations with \texttt{atomicAdd(counter, C)} until the list is exhausted.

\subsection{P2: one thread group per operation}
A group of $g\in\{1,2,4,8,16,32\}$ lanes of the same warp handles one
operation. The group examines a window of $g$ consecutive slots at once: lane
$l$ reads $T[s+l]$, so a probe step of a group is a single coalesced access (a
32-lane group reads exactly one 128-byte line). The group first asks with
\emph{one} \texttt{\_\_ballot\_sync} whether any lane saw $k$, \texttt{EMPTY} or
\texttt{TOMB}. In a full table most windows contain only other keys, so the
detailed ballots (\texttt{v==k}: found; \texttt{v==EMPTY}: end of chain;
\texttt{v==TOMB}: reusable slot) run only for the few windows where something
happens. This single-ballot test made P2 9--17\% faster than our first version,
which always issued two ballots per window. For an insert the lowest empty lane
(\texttt{\_\_ffs}) becomes the leader, performs the CAS and broadcasts the old
value with \texttt{\_\_shfl\_sync}. If the CAS fails because another key took
the slot, the next empty lane of the same window is tried, otherwise the window
advances by $g$. Tombstones are tracked exactly as in P1 (the first one on the
path, found with a ballot and broadcast by \texttt{\_\_shfl\_sync}).
\texttt{gmask} contains only the lanes of the group, so different groups of a
warp can be at different operations (independent thread scheduling on Volta+);
all lanes of a group always execute the same operations in the same order, so
every \texttt{\_sync} intrinsic is reached by the whole group.

Operations are assigned to groups in one of two ways. \textbf{Static}
(cyclic, the default): group $q$ takes operations $q, q+Q, \dots$ ($Q$ =
number of groups). \textbf{Dynamic}: the first lane of the group takes a chunk
of $C$ operations with \texttt{atomicAdd} on a per-kernel counter (one atomic per
group) and broadcasts the start to the other lanes with \texttt{\_\_shfl\_sync}.

\textbf{How these P2 improvements were found.} Both came out of our correctness
work rather than from planned tuning. When the stricter verification
exposed the tombstone bug, the P2 insert had to be restructured to track
tombstones inside a window. Adding a third ballot (\texttt{v==TOMB}) to every
window would have made every probe step more expensive, so we asked instead
whether a window needs any detailed ballot at all. In a full table almost every
window holds only other keys, so a single combined ballot answers that, and the
detailed ballots run only rarely. We measured the old and the new P2 on the same
workloads before accepting the change. P1 was unchanged (within 2\%), while P2 was
9--17\% faster, so all P2 experiments were re-run with the new version. Re-examining
P2's results for that re-run also made us question why static P2 needs
$2^{26}$ threads (16 times more threads than operations, so only about two
operations per group) to perform well. Our
guess was load imbalance caused by the expensive late inserts, which led us to try
dynamic group assignment: it removes that imbalance with a shared counter instead
of relying on the block scheduler.

\subsection{Results}\label{sec:q1res}
All tuning runs use $2^{22}$ operations with 100\% inserts (the most
expensive mix; all mixes are at least 50\% inserts).

\IfFileExists{q1_results.tex}{\subsubsection*{P1: static vs.\ dynamic assignment}
\inputif{q1_p1_static.tex}
\inputif{q1_p1_cyclic.tex}
\inputif{q1_p1_chunk.tex}
\inputif{q1_p1_dyn.tex}
\inputif{q1_p1_sched.tex}

\begin{itemize}
\item \textbf{Static and cyclic assignment depend strongly on the thread count}
  (Tables~\ref{tab:p1static}, \ref{tab:p1cyclic}): from $2^{16}$ to $2^{22}$
  threads the time drops from 9--14~s to about 1.5~s (3 repetitions per configuration). With few threads, each thread
  owns a fixed list of many operations that it must finish alone. The cost of an
  operation varies a lot (a few probes early on, thousands of probes
  when the table is nearly full), so the threads that happen to own expensive
  operations decide the run time while the others have finished and sit idle (load
  imbalance). Static assignment is only good when every thread gets a single
  operation ($2^{22}$ threads $=$ number of operations). Then the hardware
  block scheduler does the balancing: as soon as a block finishes, a new one
  starts in its place.
\item \textbf{Dynamic assignment with small chunks is insensitive to the thread
  count} (Table~\ref{tab:p1dyn}: 1.37--1.53~s for every configuration from
  $2^{16}$ to $2^{22}$ threads), because a thread that finishes early simply takes
  more work. Even $2^{16}$ threads (only 2.7$\times$ the
  \devmaxresidentthreads{} resident threads) are enough.
\item \textbf{The optimal chunk size is 1} (Table~\ref{tab:p1chunk}), and the time
  rises steadily with the chunk size (about 1.4~s at chunk 1, about 17~s at chunk 256).
  We expected larger chunks to help, as they do on the CPU, where they amortise the
  cost of the shared counter, but here the opposite happened. We think there are two
  reasons. (i) The \texttt{atomicAdd} on one counter is
  cheap compared with a single operation here (hundreds to thousands of
  dependent memory accesses), so there is nothing to amortise. (ii) Operations become more expensive the
  fuller the table is, so the last chunks handed out are the most
  expensive ones. A large last chunk makes one thread execute many of the most expensive
  operations \emph{sequentially} at the end, while the rest of the GPU is idle. Chunk 1
  minimises this tail.
\item \textbf{Block size} has only a second-order effect for P1 (Table~\ref{tab:p1dyn}):
  a thread does not cooperate with the other threads of its block, so the block
  size only affects how many blocks fit on an SM. 1024 was marginally best and
  is used below.
\item Best P1 configurations (Table~\ref{tab:p1sched}): dynamic (chunk \ponechunk)
  \ponedynamicms~ms, static \ponestaticms~ms, cyclic \ponecyclicms~ms. With one operation per
  thread the three schemes are nearly identical; dynamic wins slightly and, more
  importantly, does not need the thread count to match the operation count.
\end{itemize}

\subsubsection*{P2: thread-group size}
\inputif{q1_p2_group.tex}
\inputif{q1_p2_tune.tex}

\begin{itemize}
\item \textbf{The larger the group, the faster} (Table~\ref{tab:p2group}); $g=32$ is
  optimal at every thread count. A group of $g$ lanes examines $g$ consecutive
  slots with one coalesced load, so a probe sequence of length $L$ needs about $L/g$
  load instructions, ballots and loop iterations instead of $L$. With $g=32$ each
  probe step consumes exactly one 128-byte line, i.e.\ every byte that is
  fetched is used. Going from $g=16$ to $g=32$ still saves about a third of the time. Probe
  sequences in a full table are long (hundreds to thousands of slots), so
  even 32-slot windows are mostly useful.
\item \textbf{$g=1$ is about 2.5$\times$ slower than P1}, although it also uses one
  thread per operation. Each lane of a warp is its own group and runs its own probe loop
  containing \texttt{\_\_ballot\_sync}/\texttt{\_\_shfl\_sync} with a one-lane mask. Our
  explanation is that these synchronising intrinsics stop the compiler from running the
  32 independent loops of a warp as one predicated SIMD loop (as it does for P1), so
  the scheduler executes the lanes' loops as separate sub-warps (independent
  thread scheduling, lecture 6) and SIMD efficiency drops. We did not confirm this by looking at the generated code.
\item \textbf{Many more threads than can be resident help P2}: the best thread
  count is $2^{26}$, i.e.\ $2^{21}$ groups of 32 for $2^{22}$ operations, so each group
  gets only two operations. Like the static P1, P2 assigns operations to
  groups statically (cyclically). Our understanding is that making the work per group as
  small as possible lets the block scheduler balance the load, which removes most of
  the slow tail caused by the expensive late inserts. With $2^{20}$ threads (one group per 128 operations) P2
  is 4$\times$ slower than with $2^{26}$.
\item \textbf{Block size}: 128 is best (Table~\ref{tab:p2tune}). 1024-thread blocks
  are worse. Our explanation is that a block can only be replaced when \emph{all} of its 32 warps are
  done, so one warp stuck on a long probe sequence holds 1024 thread slots of the
  SM idle. Smaller blocks release their resources sooner. Below 128 the per-SM
  limit of \devmaxblockspersm{} blocks keeps the SM from being filled (16 blocks
  $\times$ 64 threads = 1024 $<$ \devmaxthreadspersm).
\end{itemize}

\subsubsection*{P2 with dynamic group assignment}
\inputif{q1_p2dyn_chunk.tex}
\inputif{q1_p2dyn_group.tex}

\begin{itemize}
\item \textbf{Dynamic assignment removes static P2's dependence on the thread
  count} (Table~\ref{tab:p2dynchunk}). With chunk 1 the time is 260--266~ms for
  every thread count from $2^{16}$ to $2^{24}$, whereas static P2 needs
  $2^{26}$ threads to reach 351~ms; with $2^{16}$ threads it takes 4.0~s, about
  15$\times$ longer than dynamic P2 with the same number of threads. Static P2 balanced the expensive late inserts only indirectly, by
  making the work per group so small that the block scheduler could balance it.
  A shared counter balances it directly: a group that finishes early takes the
  next operation.
\item \textbf{The best dynamic P2 (\ptwodynbestms~ms) takes 26\% less time than the best
  static P2 (\ptwobestms~ms)}, even though we tried static P2 with every thread
  count. We think the reason is that static P2 needs 524\,288 blocks to balance its
  load, and most of them start, find no work and exit, while dynamic P2 does the same
  balancing with a few thousand long-lived groups.
\item \textbf{Chunk 1 is optimal}, and the time grows steadily with the chunk
  size (about 260~ms at chunk 1, about 550~ms at chunk 16), for the same reason as
  in P1. One \texttt{atomicAdd} per group is negligible compared with an operation
  (tens to thousands of probes), while larger chunks put many of the most
  expensive operations, the last inserts into the full table, into one group's
  sequential tail.
\item \textbf{Group size 32 remains best} (Table~\ref{tab:p2dyngroup}). This is
  what we expected, since the coalescing advantage of wide windows should not depend
  on how operations are assigned to groups.
\end{itemize}

\subsubsection*{P1 vs.\ P2}
\inputif{q1_final.tex}

Table~\ref{tab:q1final} compares the best P1 (dynamic, chunk~\ponechunk, block size
\ponebs) with the best static P2 (group~\ptwogroup, block size \ptwobs) and the
best dynamic P2 (chunk~\ptwodynchunk, group~\ptwodyngroup), each at its best
thread count, for the five operation mixes and $2^{20}$, $2^{22}$ and $2^{24}$
operations. \qoneallverified{} runs passed verification. Dynamic P2 is the fastest
program in every case: \qonedynspeedupmin--\qonedynspeedupmax$\times$ faster than P1
and \qonedynvsstaticmin--\qonedynvsstaticmax$\times$ faster than static P2.
\begin{itemize}
\item \textbf{P2 is faster in every case, by \qonespeedupmin--\qonespeedupmax$\times$.}
  We think the difference comes mainly from memory-access coalescing and cooperation. P2 finds the target slot of an
  operation in $L/32$ steps, each a fully used 128-byte access, while a P1 thread walks the same
  $L$ slots one dependent load at a time. In addition, the 32 lanes of a P1 warp probe 32
  unrelated parts of the table, so every warp load touches up to 32 different lines
  (memory divergence: the warp waits for the slowest lane). P1's loads still
  hit in L1 for consecutive slots of the same thread, which is probably why the gap
  is ``only'' 4--6$\times$ and not 32$\times$.
\item \textbf{Time grows much faster than the operation count}: for 100\% inserts,
  $2^{20}\to2^{22}\to2^{24}$ operations take 0.15~s $\to$ 1.4~s $\to$ 67~s (P1)
  and 31~ms $\to$ 0.35~s $\to$ 14.9~s (static P2), although the operation count only
  grows by 4$\times$ each time. The table is sized to the number of distinct keys, so
  it ends 100\% full. For linear probing the total cost of filling a table
  completely grows like $n^{3/2}$: the last few percent of inserts have to walk
  very large clusters that keep merging. At $2^{24}$ the growth is even steeper,
  and we think the effect described next is part of the reason.
\item \textbf{Modulo bias of $h(k)=k \bmod n$ seems to matter in a full table.} Keys are uniform in
  $[0,N)$ with $N=2^{30}$. If $N/n$ is not an integer, $k \bmod n$ is not uniform: with
  $N/n = q + f$ the first $f\cdot n$ slots receive keys with probability proportional to
  $q+1$, the rest proportional to $q$. In a half-empty table such a small imbalance
  ($1/q$) probably does not matter. In a 100\% full table, however, every surplus key of the
  over-loaded region must overflow into the under-loaded one, which we expected to
  create one very large cluster that all late inserts walk through. We tested this
  with P2 on $2^{23}$ distinct insert keys:
  \begin{center}\small
  \begin{tabular}{rrr}\toprule
  table size $n$ & $N/n$ & time (ms)\\\midrule
  8\,388\,608 ($=2^{23}$) & 128.000 & 667\\
  8\,388\,000 & 128.009 & 695\\
  8\,356\,092 & 128.5 & 2374\\\bottomrule
  \end{tabular}
  \end{center}
  The same keys run $3.5\times$ slower when $N/n$ has fractional part $\frac12$ (half the
  table over-loaded by $1/128$). For our workloads $N/n\approx1024.5$, $256.5$ and $64.5$
  for $2^{20}$, $2^{22}$ and $2^{24}$ operations: the relative imbalance $1/q$ grows
  from 0.1\% to 1.6\%, which we think explains at least part of the extra jump at
  $2^{24}$.
\item \textbf{Effect of the operation mix.} Most of the cost comes from the inserts,
  because they fill the table, but the deletes and searches are not free either:
  \begin{itemize}
  \item 100/0/0 (all inserts) is the most expensive mix at $2^{22}$ and $2^{24}$, but
    not at $2^{20}$, where the two 80\% mixes take longer for all three programs
    (28--30\% for P1, 9--18\% for static P2 and 65--66\% for dynamic P2). We first thought this was due to the particular key sets, so we
    generated eight random key sets for each of these mixes at $2^{20}$ (P2 dynamic,
    $2^{20}$ threads). With every one of the eight seeds both 80\% mixes were slower than
    100/0/0 (25.9--31.0~ms against 23.3--28.9~ms), so the effect is real. Timing the insert phase alone (by giving the
    programs only the insert count), the inserts of an 80\% list take 17.6--19.3~ms
    against 25.2--28.0~ms for a 100\% list, a ratio of about 0.7, close to the
    $0.8^{3/2}\approx0.72$ that the $n^{3/2}$ growth predicts. The extra 11--12~ms
    come from the 20\% deletes or searches: in a full table each of them walks from
    the key's home slot to the key, which is a long way on average. We measured the
    same split at the larger sizes (P2 dynamic, best thread count). The deletes or
    searches cost 58--69\%, 41--49\% and 26--27\% of the time of the insert phase
    at $2^{20}$, $2^{22}$ and $2^{24}$, so their share shrinks as the table grows. At
    the same time the inserts of an 80\% list take only 0.43 ($2^{22}$) and 0.51
    ($2^{24}$) of the time of a 100\% list, much less than the 0.72 of the $n^{3/2}$
    rule. All these lists have $N/n$ with a fractional part of about 0.5, so we think
    the reason is the modulo bias described above: the large cluster it creates grows
    quickly with $n$ and makes the larger 100\% table relatively more expensive. So
    at $2^{22}$ and $2^{24}$ the extra inserts of the 100\% list cost more than the
    deletes or searches of the 80\% lists, and 100/0/0 is the most expensive mix.
  \item 50/50/0 is the cheapest at $2^{24}$, which surprised us at first. Every delete
    needs its own distinct key, so in this mix all inserted keys are distinct and
    $n$ is \emph{exactly} $2^{23}$. Then $N/n=128$ is an integer, there is no modulo
    bias, and the insert phase costs only 0.67~s instead of 2.39~s for the equally
    long insert phase of 50/0/50, whose $n$ is slightly smaller than $2^{23}$
    because of duplicate keys (we timed the insert phases separately). Each
    delete is a positive search plus one CAS.
  \item 50/0/50 and 50/50/0 are otherwise similar (at $2^{20}$ and $2^{22}$, where
    neither $n$ is a power of two): both do the same inserts, and a positive search
    and a delete both walk to the key and stop.
  \item 80/20/0 and 80/0/20 cost almost the same, for the same reason; a delete
    adds only one CAS.
  \end{itemize}
\item \textbf{The speedup of P2 is somewhat smaller at $2^{24}$ (3.9--4.5$\times$) than at
  $2^{20}$ (4.8--5.8$\times$).} We are not completely sure why. One possible
  explanation is that at the largest size both programs spend most of their time
  walking very long clusters near the end, where P1's sequential walk also uses whole
  cache lines through L1, and both programs compete for the few remaining empty
  slots.
\end{itemize}

\subsubsection*{Robustness across random key sets}
\inputif{q1_seeds.tex}

To check that the conclusions do not depend on one particular random operation
list, we regenerated the $2^{22}$-operation lists of all five mixes with two
further seeds and ran the best configuration of each program on them
(Table~\ref{tab:q1seeds}).
\begin{itemize}
\item \textbf{Absolute times vary by up to about 50\% between key sets} of the same
  mix (6--47\% depending on mix and program, e.g.\ 50/50/0: 871--1181~ms for P1). Our first guess was the modulo bias
  above, but the data does not support it. For every mix the three seeds have almost the
  same $N/n$ (fractional part $\approx0.49$), and in the 50/50/0 mix $n=2^{21}$
  exactly, so there is no bias at all, yet that mix varies by 35\%. We think the
  variation comes from the 100\% full linearly probed table itself: its cost is
  dominated by the few largest clusters, which the last inserts must walk, and the
  size of those clusters depends strongly on the exact keys (the cost seems to be
  heavy-tailed).
\item \textbf{The comparisons are stable}: P1/P2 lies in
  \qoneseedratioonemin--\qoneseedratioonemax$\times$ and static/dynamic P2 in
  \qoneseedratiotwomin--\qoneseedratiotwomax$\times$ across all seeds, and the
  ranking dynamic P2 $<$ static P2 $<$ P1 holds in all \qoneseedsameorder{} of
  the 15 seed/mix combinations. So we think the conclusions of this section do not
  depend on the particular random lists. Absolute times of a single list should
  be read with this key-set variation (up to about 50\%) in mind; the repetitions
  (3--5 per point) only remove timing noise, not this variation.
\end{itemize}

}{}

%=====================================================================
\section{Problem 2: $x^TAy$}

\subsection{Programs}
$A$ is stored row-major; element $e$ of the flattened matrix has
$i=e\gg\log_2 n$ and $j=e\,\&\,(n-1)$. A separate \texttt{init\_kernel}
fills $A$, $x$, $y$ with values from a hash of the element index (exact
multiples of $2^{-15}$ in $[-1,1)$), so the CPU can regenerate the same
floats and compute the reference $x^TAy$ in double precision without copying
$A$ back. Initialisation, the CPU reference and the difference are not
timed. All arithmetic on the GPU is in \texttt{float}.

\begin{enumerate}
\item \textbf{Basic} (\texttt{q2\_v1\_basic.cu}): thread $id$ handles the
  contiguous chunk of $mn/t$ elements starting at $id\cdot mn/t$ and
  accumulates $x_iA_{ij}y_j$; partial sums are written to an array, copied to
  the host and added there. Within a warp, the 32 threads access addresses
  $4mn/t$ bytes apart: every load instruction of a warp touches 32 different
  cache lines.
\item \textbf{Coalesced} (\texttt{q2\_v2}): thread $id$ handles $id, id+t,
  id+2t,\dots$; in every iteration a warp reads 32 consecutive floats of $A$
  (one 128-byte line) and of $y$. Host merge as before.
\item \textbf{Tree reduction} (\texttt{q2\_v3}): coalesced accumulation, then
  a binary-tree reduction of the block's partial sums ($\log_2$ blocksize
  halving steps separated by \texttt{\_\_syncthreads}) in a global-memory
  scratch array; thread 0 adds the block total to the result with one
  \texttt{atomicAdd}. Only 4 bytes are copied back.
\item \textbf{Shared memory} (\texttt{q2\_v4}): same tree, but in a
  \texttt{\_\_shared\_\_} array.
\item \textbf{Additional optimisations} (\texttt{q2\_v5}): \texttt{float4}
  loads of $A$ and $y$ (one 16-byte load per thread, 512 bytes per warp
  instruction, $4\times$ fewer instructions/iterations), read-only cache path
  (\texttt{\_\_ldg}, \texttt{const \_\_restrict\_\_}), loop unrolled by 2 with
  two independent accumulators (two loads in flight per thread), and a
  reduction with \texttt{\_\_shfl\_down\_sync} inside each warp, one
  shared-memory slot per warp, and a final warp-shuffle reduction.
\end{enumerate}
The OpenMP program (\texttt{q2\_omp.c}) splits rows statically among threads,
computes each row dot product $\sum_j A_{ij}y_j$ with an
\texttt{omp simd reduction} loop (AVX2) and accumulates $x_i\cdot$row in a
double \texttt{reduction(+)}; $A$ is first-touched by the same threads.

\subsection{Results}
The largest square size that fits comfortably in the 4~GB GPU is
$m=n=2^{14}$ ($A$ is 1~GB; $2^{15}$ would need 4~GB). $x^TAy$ reads every
element of $A$ exactly once and does only 2--3 flops per 4 bytes, so we expected
it to be limited by memory bandwidth, and we report the achieved bandwidth
($4(mn+m+n)$ bytes / time) relative to the \devpeakbandwidthgbs~GB/s peak.

\IfFileExists{q2_results.tex}{\subsubsection*{Basic program: thread count and thread hierarchy}
\inputif{q2_basic.tex}
\inputif{q2_basic_bs.tex}

Tables~\ref{tab:q2basic} and~\ref{tab:q2basicbs} show the basic program for
$m=n=2^{14}$. The best configuration is $t=$\,\qtwobasicbestt{} threads in blocks of
\qtwobasicbestbs{} (\qtwobasicbestms~ms).
\begin{itemize}
\item \textbf{Up to $2^{10}$ threads} the time halves every time $t$ doubles, and it
  flattens out between $2^{11}$ and $2^{12}$.
  There are too few threads to hide the DRAM latency: with $t\le 2^{10}$ there
  are at most 32 warps on 16 SMs, so most SMs have one or two warps, and each
  warp stalls for hundreds of cycles on every load.
\item \textbf{$2^{12}$--$2^{13}$ threads}: the time flattens at about 16.6--17~ms,
  i.e.\ roughly 64~GB/s, which is only about 37\% of peak. A warp load touches 32 different
  cache lines (one per thread, $4mn/t$ bytes apart). We think it reaches even
  this only because each thread walks sequentially through its chunk, so after the first
  miss the next 31 loads of the same thread hit the same 128-byte line in L1.
\item \textbf{Block size matters a lot beyond $2^{13}$ threads}. With blocksize 32 the time stays at 16--17~ms
  up to $2^{17}$ threads, while every larger block size jumps to 26--32~ms.
  Our explanation is the L1 working set (we did not measure cache hit rates
  with a profiler, so this is only an estimate from the cache sizes). Each resident warp keeps 32 lines
  (4~KB) live in L1. An SM can hold 48 warps (1536 threads), which is a 192~KB
  working set, larger than the 128~KB L1/shared array, so lines are evicted
  before their 32 elements are consumed and are fetched again. With 32-thread
  blocks the limit of 16 blocks per SM (compute capability 8.6) caps
  residency at 16 warps per SM (64~KB), which fits in L1. So here lower occupancy
  seems to help, because it avoids cache thrashing. For $t\le 2^{13}$ there are
  at most 16 warps per SM anyway, so there the block size matters only when there are
  fewer blocks than SMs: with $t=2^{13}$ and 1024-thread blocks, or $t=2^{12}$ and
  512-thread blocks, only 8 blocks exist, so only 8 of the 16 SMs work (30.9 and
  31.1~ms instead of about 17~ms).
\item \textbf{Very large $t$} ($\ge 2^{19}$) is slower again. The chunks get short
  (512 elements and less), and the program has to copy $t$ partial
  sums to the host and add them there. Table~\ref{tab:q2trend} shows this cost:
  versions 2 and 3 differ only in the merge, and the gap grows with $t$.
\end{itemize}

\subsubsection*{Effect of the optimisations}
\inputif{q2_versions.tex}
\inputif{q2_trend.tex}

\begin{itemize}
\item \textbf{Coalescing} is the big step: $2.5\times$ over the best basic
  configuration (16.6~ms $\to$ 6.6~ms). Each warp load
  is now one 128-byte transaction that is fully used, so the program reaches
  about 160~GB/s, which is 92\% of peak. Coalescing needs enough threads
  (Table~\ref{tab:q2trend}). With $t\le 2^{12}$ the coalesced version is
  \emph{slower} than the basic one, because each thread then issues one independent
  load per iteration and there are not enough warps to keep enough requests in
  flight to cover the DRAM latency. The basic version is helped there by its L1 hits.
\item \textbf{Tree reduction} on the GPU removes the copy of $t$ partial sums
  and the serial host merge. At the thread count where coalescing is best the gain is
  small (the merge of $2^{16}$ floats is cheap; at $2^{14}$ threads the global-memory tree
  with its $\log_2(\text{blocksize})$ barrier steps is even slightly slower than the host merge), but the gain grows quickly
  with $t$ ($2^{22}$ threads: 13.1~ms $\to$ 6.7~ms). So with the tree reduction
  the best time no longer depends much on $t$.
\item \textbf{Shared memory} for the tree changes little (1.6--6\%). The
  reduction is $\log_2(\text{blocksize})$ steps over at most 1024 values once
  per block, which is negligible compared with streaming 1~GB from DRAM. Its
  global-memory version mostly hits in L1/L2 anyway. Shared memory helps most
  at large $t$, where there are many blocks and therefore many reductions.
\item \textbf{float4 + ILP + warp shuffle} gives the best time and,
  more importantly, reaches it at only $2^{12}$--$2^{13}$ threads. Each load
  instruction moves 16 bytes per thread, and two independent loads per thread are in
  flight, so fewer warps can still keep the memory system busy. The
  shuffle reduction needs no shared memory and only one \texttt{\_\_syncthreads}.
\end{itemize}
All optimised versions converge to 6.3--6.6~ms $\approx$ 160--171~GB/s; the best (v5,
\qtwovfivems~ms) reaches \qtwovfivegbps~GB/s, i.e.\ 97\% of the \devpeakbandwidthgbs~GB/s peak. Since
the kernel is limited by memory bandwidth, we do not expect further optimisations
to gain much; we think the remaining few percent are lost to things like DRAM
refresh and the start and end of the kernel. The numerical difference from the double-precision CPU result is about
$10^{-4}$--$10^{-3}$ for a result of magnitude $\approx 3300$ (relative error
$10^{-7}$--$10^{-6}$). The difference changes between versions because the order of
the float additions changes; it looks like versions that sum many short partial sums
(large $t$) are more accurate than those with long sequential float accumulations
(small $t$).

\subsubsection*{Best CUDA program vs.\ OpenMP}
\inputif{q2_compare.tex}
\inputif{q2_omp.tex}

The best CUDA program (v5, tuned per size) is \qtwospeedupmin--\qtwospeedupmax$\times$
faster than the best OpenMP program (Table~\ref{tab:q2compare}). For every matrix of
64~MB or more the speedup is about $4.3$--$4.7\times$, and it is the same for
square, wide ($2^{12}\times2^{16}$) and tall ($2^{16}\times2^{12}$)
matrices. Both programs seem to be limited by memory bandwidth, so the speedup is roughly
the ratio of the achieved memory bandwidths: about 170~GB/s on the GPU's GDDR6 against 35--40~GB/s
on the CPU's DDR4 (Table~\ref{tab:q2omp}). The CPU seems to reach its memory
bandwidth limit with only 2--4 threads (2 and 4 threads give almost the same time); more threads only seem to add contention and
synchronisation (the 16-thread runs are slower, and for the smallest size the
fork/join cost of 16 threads dominates). For the smallest matrix
($2^{10}\times2^{10}$, 4~MB) there is no gain (both take about 45~$\mu$s): the run time is
tens of microseconds, so we think kernel-launch and copy latency dominate, while the
CPU's 4~MB working set stays in its 16~MB L3 cache between repetitions.
}{}

\end{document}
